\documentclass{amsart}
% For dealing with references we use the comment environment
\usepackage{verbatim}
\newenvironment{reference}{\comment}{\endcomment}
%\newenvironment{reference}{}{}
\newenvironment{slogan}{\comment}{\endcomment}
\newenvironment{history}{\comment}{\endcomment}

% For commutative diagrams we use Xy-pic
\usepackage[all]{xy}

% We use 2cell for 2-commutative diagrams.
\xyoption{2cell}
\UseAllTwocells

% We use multicol for the list of chapters between chapters
\usepackage{multicol}

% This is generally recommended for better output
\usepackage{lmodern}
\usepackage[T1]{fontenc}

% For cross-file-references

% Package for hypertext links:
\usepackage{hyperref}

% For any local file, say "hello.tex" you want to link to please

% Theorem environments.
%
\theoremstyle{plain}
\newtheorem{theorem}[subsection]{Theorem}
\newtheorem{proposition}[subsection]{Proposition}
\newtheorem{lemma}[subsection]{Lemma}

\theoremstyle{definition}
\newtheorem{definition}[subsection]{Definition}
\newtheorem{example}[subsection]{Example}
\newtheorem{exercise}[subsection]{Exercise}
\newtheorem{situation}[subsection]{Situation}

\theoremstyle{remark}
\newtheorem{remark}[subsection]{Remark}
\newtheorem{remarks}[subsection]{Remarks}

\numberwithin{equation}{subsection}

% Macros
%
\def\lim{\mathop{\mathrm{lim}}\nolimits}
\def\colim{\mathop{\mathrm{colim}}\nolimits}
\def\Spec{\mathop{\mathrm{Spec}}}
\def\Hom{\mathop{\mathrm{Hom}}\nolimits}
\def\Ext{\mathop{\mathrm{Ext}}\nolimits}
\def\SheafHom{\mathop{\mathcal{H}\!\mathit{om}}\nolimits}
\def\SheafExt{\mathop{\mathcal{E}\!\mathit{xt}}\nolimits}
\def\Sch{\mathit{Sch}}
\def\Mor{\mathop{\mathrm{Mor}}\nolimits}
\def\Ob{\mathop{\mathrm{Ob}}\nolimits}
\def\Sh{\mathop{\mathit{Sh}}\nolimits}
\def\NL{\mathop{N\!L}\nolimits}
\def\CH{\mathop{\mathrm{CH}}\nolimits}
\def\proetale{{pro\text{-}\acute{e}tale}}
\def\etale{{\acute{e}tale}}
\def\QCoh{\mathit{QCoh}}
\def\Ker{\mathop{\mathrm{Ker}}}
\def\Im{\mathop{\mathrm{Im}}}
\def\Coker{\mathop{\mathrm{Coker}}}
\def\Coim{\mathop{\mathrm{Coim}}}

% Boxtimes
%
\DeclareMathSymbol{\boxtimes}{\mathbin}{AMSa}{"02}

%
% Macros for moduli stacks/spaces
%
\def\QCohstack{\mathcal{QC}\!\mathit{oh}}
\def\Cohstack{\mathcal{C}\!\mathit{oh}}
\def\Spacesstack{\mathcal{S}\!\mathit{paces}}
\def\Quotfunctor{\mathrm{Quot}}
\def\Hilbfunctor{\mathrm{Hilb}}
\def\Curvesstack{\mathcal{C}\!\mathit{urves}}
\def\Polarizedstack{\mathcal{P}\!\mathit{olarized}}
\def\Complexesstack{\mathcal{C}\!\mathit{omplexes}}
% \Pic is the operator that assigns to X its picard group, usage \Pic(X)
% \Picardstack_{X/B} denotes the Picard stack of X over B
% \Picardfunctor_{X/B} denotes the Picard functor of X over B
\def\Pic{\mathop{\mathrm{Pic}}\nolimits}
\def\Picardstack{\mathcal{P}\!\mathit{ic}}
\def\Picardfunctor{\mathrm{Pic}}
\def\Deformationcategory{\mathcal{D}\!\mathit{ef}}

\setlength{\emergencystretch}{3em}
\begin{document}
\title[Stacks Project, Tag 08Z9]{1747 --- Matlis duality exchanges Artinian and Noetherian modules over a complete local Noetherian ring via the residue-field injective hull}
\maketitle
\noindent Copyright (C) 2005--2025 Johan de Jong. Stacks Project authors. Source: \href{https://stacks.math.columbia.edu/tag/08Z9}{Tag 08Z9}.

\noindent Editorial difficulty note (not author text). Difficulty H2: Matlis duality exchanges all Artinian and Noetherian modules over a complete local Noetherian ring, including modules that are not finite length. Its retained proof core reconstructs the completion as the endomorphism ring of the injective hull, derives the hull descending-chain condition from Noetherian completion, and builds finite injective copresentations from finite socles and essential torsion embeddings before transporting presentations through an exact contravariant functor. This completion/injective-hull structural duality is substantially beyond finite-dimensional vector-space duality or ordinary qualifications..

\noindent Editorial provenance note (not author text). History: excerpt from revision \texttt{a0444\allowbreak 6e57e\allowbreak c1fbc\allowbreak 252a8\allowbreak 71afc\allowbreak ec775\allowbreak 2fb28\allowbreak 07b14}, dualizing.tex, lines 967--996. Reference presentation was adapted; original-author context and core support blocks were retained or added. The mathematical wording is unchanged. Licensed under GNU FDL 1.2 or later; no invariant sections or cover texts. See ../licenses/COPYING. Context blocks reproduce author definitions and supporting results; they are not additional counted problems.

\begin{definition}
\label{definition-essential}
Let $\mathcal{A}$ be an abelian category.
\begin{enumerate}
\item An injection $A \subset B$ of $\mathcal{A}$ is {\it essential},
or we say that $B$ is an {\it essential extension of} $A$,
if every nonzero subobject $B' \subset B$ has nonzero intersection with $A$.
\item A surjection $f : A \to B$ of $\mathcal{A}$ is {\it essential}
if for every proper subobject $A' \subset A$ we have $f(A') \not = B$.
\end{enumerate}
\end{definition}

\begin{definition}
\label{definition-injective-hull}
Let $R$ be a ring. An injection $M \to I$ of $R$-modules is said
to be an {\it injective hull} if $I$ is an injective $R$-module and
$M \to I$ is an essential injection.
\end{definition}

\begin{lemma}
\label{lemma-endos}
Let $(R, \mathfrak m, \kappa)$ be a Noetherian local ring.
Let $E$ be an injective hull of $\kappa$ over $R$. Then
$\Hom_R(E, E)$ is canonically isomorphic to the completion of $R$.
\end{lemma}

\begin{proof}
Write $E = \bigcup E_n$ with $E_n = E[\mathfrak m^n]$ as in
Lemma \ref{lemma-union-artinian}. Any endomorphism of $E$
preserves this filtration. Hence
$$
\Hom_R(E, E) = \lim \Hom_R(E_n, E_n)
$$
The lemma follows as
$\Hom_R(E_n, E_n) = \Hom_{R/\mathfrak m^n}(E_n, E_n) = R/\mathfrak m^n$
by Lemma \ref{lemma-evaluate}.
\end{proof}


\begin{lemma}
\label{lemma-describe-categories}
Let $(R, \mathfrak m, \kappa)$ be a Noetherian local ring.
Let $E$ be an injective hull of $\kappa$.
\begin{enumerate}
\item For an $R$-module $M$ the following are equivalent:
\begin{enumerate}
\item $M$ satisfies the ascending chain condition,
\item $M$ is a finite $R$-module, and
\item there exist $n, m$ and an exact sequence
$R^{\oplus m} \to R^{\oplus n} \to M \to 0$.
\end{enumerate}
\item For an $R$-module $M$ the following are equivalent:
\begin{enumerate}
\item $M$ satisfies the descending chain condition,
\item $M$ is $\mathfrak m$-power torsion and
$\dim_\kappa(M[\mathfrak m]) < \infty$, and
\item there exist $n, m$ and an exact sequence
$0 \to M \to E^{\oplus n} \to E^{\oplus m}$.
\end{enumerate}
\end{enumerate}
\end{lemma}

\begin{proof}
We omit the proof of (1).

\medskip\noindent
Let $M$ be an $R$-module with the descending chain condition. Let $x \in M$.
Then $\mathfrak m^n x$ is a descending chain of submodules, hence stabilizes.
Thus $\mathfrak m^nx = \mathfrak m^{n + 1}x$ for some $n$. By Nakayama's lemma
(Algebra, Lemma \href{https://stacks.math.columbia.edu/tag/00DV}{lemma NAK (Tag 00DV)}) this implies $\mathfrak m^n x = 0$,
i.e., $x$ is $\mathfrak m$-power torsion. Since $M[\mathfrak m]$ is a vector
space over $\kappa$ it has to be finite dimensional in order to have the
descending chain condition.

\medskip\noindent
Assume that $M$ is $\mathfrak m$-power torsion and has a finite dimensional
$\mathfrak m$-torsion submodule $M[\mathfrak m]$. By
Lemma \ref{lemma-torsion-submodule-sum-injective-hulls}
we see that $M$ is a submodule of $E^{\oplus n}$ for some $n$.
Consider the quotient $N = E^{\oplus n}/M$. By
Lemma \ref{lemma-injective-hull-has-dcc} the module $E$ has the
descending chain condition hence so do $E^{\oplus n}$ and $N$.
Therefore $N$ satisfies (2)(a) which implies $N$ satisfies
(2)(b) by the second paragraph of the proof. Thus by
Lemma \ref{lemma-torsion-submodule-sum-injective-hulls}
again we see that $N$ is a submodule of $E^{\oplus m}$ for some $m$.
Thus we have a short exact sequence
$0 \to M \to E^{\oplus n} \to E^{\oplus m}$.

\medskip\noindent
Assume we have a short exact sequence
$0 \to M \to E^{\oplus n} \to E^{\oplus m}$.
Since $E$ satisfies the descending chain condition by
Lemma \ref{lemma-injective-hull-has-dcc}
so does $M$.
\end{proof}


\begin{lemma}
\label{lemma-injective-hull-has-dcc}
Let $(R, \mathfrak m, \kappa)$ be a Noetherian local ring.
Let $E$ be an injective hull of $\kappa$ over $R$. Then
$E$ satisfies the descending chain condition.
\end{lemma}

\begin{proof}
If $E \supset M_1 \supset M_2 \supset \ldots$ is a sequence of submodules, then
$$
\Hom_R(E, E) \to \Hom_R(M_1, E) \to \Hom_R(M_2, E) \to \ldots
$$
is a sequence of surjections. By Lemma \ref{lemma-endos} each of these is a
module over the completion $R^\wedge = \Hom_R(E, E)$.
Since $R^\wedge$ is Noetherian
(Algebra, Lemma \href{https://stacks.math.columbia.edu/tag/0316}{lemma completion Noetherian Noetherian (Tag 0316)})
the sequence stabilizes: $\Hom_R(M_n, E) = \Hom_R(M_{n + 1}, E) = \ldots$.
Since $E$ is injective, this can only happen if $\Hom_R(M_n/M_{n + 1}, E)$
is zero. However, if $M_n/M_{n + 1}$ is nonzero, then it contains a
nonzero element annihilated by $\mathfrak m$, because $E$ is
$\mathfrak m$-power torsion by Lemma \ref{lemma-union-artinian}.
In this case $M_n/M_{n + 1}$ has a nonzero map into $E$, contradicting
the assumed vanishing. This finishes the proof.
\end{proof}


\begin{lemma}
\label{lemma-evaluate}
Let $(R, \mathfrak m, \kappa)$ be an artinian local ring.
Let $E$ be an injective hull of $\kappa$.
For any finite $R$-module $M$ the evaluation map
$$
M \longrightarrow \Hom_R(\Hom_R(M, E), E)
$$
is an isomorphism. In particular $R = \Hom_R(E, E)$.
\end{lemma}

\begin{proof}
Observe that the displayed arrow is injective. Namely, if $x \in M$ is
a nonzero element, then there is a nonzero map $Rx \to \kappa$ which
we can extend to a map $\varphi : M \to E$ that doesn't vanish on $x$.
Since the source and target of the arrow have the same length by
Lemma \ref{lemma-finite}
we conclude it is an isomorphism. The final statement follows
on taking $M = R$.
\end{proof}


\begin{lemma}
\label{lemma-finite}
Let $(R, \mathfrak m, \kappa)$ be an artinian local ring.
Let $E$ be an injective hull of $\kappa$. For every finite
$R$-module $M$ we have
$$
\text{length}_R(M) = \text{length}_R(\Hom_R(M, E))
$$
In particular, the injective hull $E$ of $\kappa$ is a finite $R$-module.
\end{lemma}

\begin{proof}
Because $E$ is an essential extension of $\kappa$ we have
$\kappa = E[\mathfrak m]$ where $E[\mathfrak m]$ is the
$\mathfrak m$-torsion in $E$ (notation as in More on Algebra, Section
\href{https://stacks.math.columbia.edu/tag/0ALX}{section torsion (Tag 0ALX)}).
Hence $\Hom_R(\kappa, E) \cong \kappa$ and the equality of lengths
holds for $M = \kappa$. We prove the displayed equality of the lemma
by induction on the length of $M$. If $M$ is nonzero there exists a surjection
$M \to \kappa$ with kernel $M'$. Since the functor $M \mapsto \Hom_R(M, E)$
is exact we obtain a short exact sequence
$$
0 \to \Hom_R(\kappa, E) \to \Hom_R(M, E) \to \Hom_R(M', E) \to 0.
$$
Additivity of length for this sequence and the sequence
$0 \to M' \to M \to \kappa \to 0$ and the equality for $M'$ (induction
hypothesis) and $\kappa$ implies the equality for $M$.
The final statement of the lemma follows as $E = \Hom_R(R, E)$.
\end{proof}


\begin{lemma}
\label{lemma-union-artinian}
Let $(R, \mathfrak m, \kappa)$ be a Noetherian local ring.
Let $E$ be an injective hull of $\kappa$ over $R$.
Let $E_n$ be an injective hull of $\kappa$ over $R/\mathfrak m^n$.
Then $E = \bigcup E_n$ and $E_n = E[\mathfrak m^n]$.
\end{lemma}

\begin{proof}
We have $E_n = E[\mathfrak m^n]$ by Lemma \ref{lemma-quotient}.
We have $E = \bigcup E_n$ because $\bigcup E_n = E[\mathfrak m^\infty]$
is an injective $R$-submodule which contains $\kappa$, see
Lemma \ref{lemma-injective-module-divide}.
\end{proof}


\begin{lemma}
\label{lemma-quotient}
Let $R \to S$ be a surjective map of local rings with kernel $I$.
Let $E$ be the injective hull of the residue field of $R$ over $R$.
Then $E[I]$ is the injective hull of the residue field of $S$ over $S$.
\end{lemma}

\begin{proof}
Observe that $E[I] = \Hom_R(S, E)$ as $S = R/I$. Hence $E[I]$ is an injective
$S$-module by Lemma \href{https://stacks.math.columbia.edu/tag/08XR}{lemma hom injective (Tag 08XR)}. Since $E$ is an essential
extension of $\kappa = R/\mathfrak m_R$ it follows that $E[I]$ is an
essential extension of $\kappa$ as well. The result follows.
\end{proof}


\begin{lemma}
\label{lemma-torsion-submodule-sum-injective-hulls}
Let $(R, \mathfrak m, \kappa)$ be a local ring.
Let $E$ be the injective hull of $\kappa$.
Let $M$ be a $\mathfrak m$-power torsion $R$-module
with $n = \dim_\kappa(M[\mathfrak m]) < \infty$.
Then $M$ is isomorphic to a submodule of $E^{\oplus n}$.
\end{lemma}

\begin{proof}
Observe that $E^{\oplus n}$ is the injective hull of
$\kappa^{\oplus n} = M[\mathfrak m]$. Thus there is an $R$-module map
$M \to E^{\oplus n}$ which is injective on $M[\mathfrak m]$.
Since $M$ is $\mathfrak m$-power torsion the inclusion
$M[\mathfrak m] \subset M$ is an essential extension
(for example by Lemma \ref{lemma-essential-extension})
we conclude that the kernel of $M \to E^{\oplus n}$ is zero.
\end{proof}


\begin{lemma}
\label{lemma-injective-module-divide}
Let $R$ be a Noetherian ring. Let $I$ be an injective $R$-module.
\begin{enumerate}
\item Let $f \in R$. Then $E = \bigcup I[f^n] = I[f^\infty]$
is an injective submodule of $I$.
\item Let $J \subset R$ be an ideal. Then the $J$-power torsion
submodule $I[J^\infty]$ is an injective submodule of $I$.
\end{enumerate}
\end{lemma}

\begin{proof}
We will use Lemma \ref{lemma-essential-extensions-in-injective}
to prove (1).
Suppose that $E \subset E' \subset I$ and that $E'$ is an essential
extension of $E$. We will show that $E' = E$. If not, then we can
find $x \in E'$ and $x \not \in E$.
Let $J = \{ a \in R \mid ax \in E\}$. Since $R$ is Noetherian,
we may write $J = (g_1, \ldots, g_t)$ for some
$g_i \in R$. By definition $E$ is the set of elements of $I$ annihilated
by powers of $f$, so we may choose integers $n_i$ so that $f^{n_i}g_ix = 0$.
Set $n = \mathrm{max}\{ n_i \}$. Then $x' = f^n x$ is an element of $E'$
not in $E$ and is annihilated by $J$. Set $J' = \{ a \in R \mid ax' \in E \}$
so $J \subset J'$. Conversely, we have $a \in J'$ if and only if $ax' \in E$
if and only if $f^m a x' = 0$ for some $m \geq 0$. But then
$f^m a x' = f^{m + n} a x$ implies $ax \in E$, i.e., $a \in J$.
Hence $J = J'$. Thus $J = J' = \text{Ann}(x')$, so $Rx' \cap E  = 0$.
Hence $E'$ is not an essential extension of $E$, a contradiction.

\medskip\noindent
To prove (2) write $J = (f_1, \ldots, f_t)$. Then
$I[J^\infty]$ is equal to
$$
(\ldots((I[f_1^\infty])[f_2^\infty])\ldots)[f_t^\infty]
$$
and the result follows from (1) and induction.
\end{proof}


\begin{lemma}
\label{lemma-essential-extensions-in-injective}
Let $R$ be a ring. Let $I$ be an injective $R$-module. Let $E \subset I$
be a submodule. The following are equivalent
\begin{enumerate}
\item $E$ is injective, and
\item for all $E \subset E' \subset I$ with $E \subset E'$ essential
we have $E = E'$.
\end{enumerate}
In particular, an $R$-module is injective if and only if every essential
extension is trivial.
\end{lemma}

\begin{proof}
The final assertion follows from the first and the fact that the
category of $R$-modules has enough injectives
(More on Algebra, Section \href{https://stacks.math.columbia.edu/tag/01D8}{section injectives modules (Tag 01D8)}).

\medskip\noindent
Assume (1). Let $E \subset E' \subset I$ as in (2).
Then the map $\text{id}_E : E \to E$ can be extended
to a map $\alpha : E' \to E$. The kernel of $\alpha$ has to be
zero because it intersects $E$ trivially and $E'$ is an essential
extension. Hence $E = E'$.

\medskip\noindent
Assume (2). Let $M \subset N$ be $R$-modules and let $\varphi : M \to E$
be an $R$-module map. In order to prove (1) we have to show that
$\varphi$ extends to a morphism $N \to E$. Consider the set $\mathcal{S}$
of pairs
$(M', \varphi')$ where $M \subset M' \subset N$ and $\varphi' : M' \to E$
is an $R$-module map agreeing with $\varphi$ on $M$. We define an ordering
on $\mathcal{S}$ by the rule $(M', \varphi') \leq (M'', \varphi'')$
if and only if $M' \subset M''$ and $\varphi''|_{M'} = \varphi'$.
It is clear that we can take the maximum of a totally ordered subset
of $\mathcal{S}$. Hence by Zorn's lemma we may assume $(M, \varphi)$
is a maximal element.

\medskip\noindent
Choose an extension $\psi : N \to I$ of $\varphi$ composed
with the inclusion $E \to I$. This is possible as $I$ is injective.
If $\psi(N) \subset E$, then $\psi$ is the desired extension.
If $\psi(N)$ is not contained in $E$, then by (2) the inclusion
$E \subset E + \psi(N)$ is not essential. hence
we can find a nonzero submodule $K \subset E + \psi(N)$ meeting $E$ in $0$.
This means that $M' = \psi^{-1}(E + K)$ strictly contains $M$.
Thus we can extend $\varphi$ to $M'$ using
$$
M' \xrightarrow{\psi|_{M'}} E + K \to (E + K)/K = E
$$
This contradicts the maximality of $(M, \varphi)$.
\end{proof}


\begin{lemma}
\label{lemma-essential-extension}
Let $R$ be a ring. Let $M \subset N$ be $R$-modules. The following
are equivalent
\begin{enumerate}
\item $M \subset N$ is an essential extension,
\item for all $x \in N$ nonzero there exists an $f \in R$ such that $fx \in M$
and $fx \not = 0$.
\end{enumerate}
\end{lemma}

\begin{proof}
Assume (1) and let $x \in N$ be a nonzero element. By (1) we have
$Rx \cap M \not = 0$. This implies (2).

\medskip\noindent
Assume (2). Let $N' \subset N$ be a nonzero submodule. Pick $x \in N'$
nonzero. By (2) we can find $f \in R$ with $fx \in M$ and $fx \not = 0$.
Thus $N' \cap M \not = 0$.
\end{proof}


\begin{lemma}
\label{lemma-injective-hull}
Let $R$ be a ring. Any $R$-module has an injective hull.
\end{lemma}

\begin{proof}
Let $M$ be an $R$-module. By
More on Algebra, Section \href{https://stacks.math.columbia.edu/tag/01D8}{section injectives modules (Tag 01D8)}
the category of $R$-modules has enough injectives.
Choose an injection $M \to I$ with $I$ an injective $R$-module.
Consider the set $\mathcal{S}$ of submodules $M \subset E \subset I$
such that $E$ is an essential extension of $M$. We order $\mathcal{S}$
by inclusion. If $\{E_\alpha\}$ is a totally ordered subset
of $\mathcal{S}$, then $\bigcup E_\alpha$ is an essential extension of $M$
too (Lemma \ref{lemma-union-essential-extensions}).
Thus we can apply Zorn's lemma and find a maximal element
$E \in \mathcal{S}$. We claim $M \subset E$ is an injective hull, i.e.,
$E$ is an injective $R$-module. This follows from
Lemma \ref{lemma-essential-extensions-in-injective}.
\end{proof}


\begin{lemma}
\label{lemma-injective-hull-unique}
Let $R$ be a ring. Let $M$, $N$ be $R$-modules and let $M \to E$
and $N \to E'$ be injective hulls. Then
\begin{enumerate}
\item for any $R$-module map $\varphi : M \to N$ there exists an
$R$-module map $\psi : E \to E'$ such that
$$
\xymatrix{
M \ar[r] \ar[d]_\varphi & E \ar[d]^\psi \\
N \ar[r] & E'
}
$$
commutes,
\item if $\varphi$ is injective, then $\psi$ is injective,
\item if $\varphi$ is an essential injection, then $\psi$ is an isomorphism,
\item if $\varphi$ is an isomorphism, then $\psi$ is an isomorphism,
\item if $M \to I$ is an embedding of $M$ into an injective $R$-module,
then there is an isomorphism $I \cong E \oplus I'$ compatible with
the embeddings of $M$,
\end{enumerate}
In particular, the injective hull $E$ of $M$ is unique up to isomorphism.
\end{lemma}

\begin{proof}
Part (1) follows from the fact that $E'$ is an injective $R$-module.
Part (2) follows as $\Ker(\psi) \cap M = 0$
and $E$ is an essential extension of $M$.
Assume $\varphi$ is an essential injection. Then
$E \cong \psi(E) \subset E'$ by (2) which implies
$E' = \psi(E) \oplus E''$ because $E$ is injective.
Since $E'$ is an essential extension of
$M$ (Lemma \href{https://stacks.math.columbia.edu/tag/08XK}{lemma essential (Tag 08XK)}) we get $E'' = 0$.
Part (4) is a special case of (3).
Assume $M \to I$ as in (5).
Choose a map $\alpha : E \to I$ extending the map $M \to I$.
Arguing as before we see that $\alpha$ is injective.
Thus as before $\alpha(E)$ splits off from $I$.
This proves (5).
\end{proof}


\begin{lemma}
\label{lemma-union-essential-extensions}
Let $R$ be a ring. Let $M$ be an $R$-module. Let $E = \colim E_i$
be a filtered colimit of $R$-modules. Suppose given a compatible
system of essential injections $M \to E_i$ of $R$-modules.
Then $M \to E$ is an essential injection.
\end{lemma}

\begin{proof}
Immediate from the definitions and the fact that filtered
colimits are exact (Algebra, Lemma \href{https://stacks.math.columbia.edu/tag/00DB}{lemma directed colimit exact (Tag 00DB)}).
\end{proof}


\begin{proposition}[Matlis duality]
\label{proposition-matlis}
Let $(R, \mathfrak m, \kappa)$ be a complete local Noetherian ring.
Let $E$ be an injective hull of $\kappa$ over $R$. The functor
$D(-) = \Hom_R(-, E)$ induces an anti-equivalence
$$
\left\{
\begin{matrix}
R\text{-modules with the} \\
\text{descending chain condition}
\end{matrix}
\right\}
\longleftrightarrow
\left\{
\begin{matrix}
R\text{-modules with the} \\
\text{ascending chain condition}
\end{matrix}
\right\}
$$
and we have $D \circ D = \text{id}$ on either side of the equivalence.
\end{proposition}

\begin{proof}
By Lemma \ref{lemma-endos} we have $R = \Hom_R(E, E) = D(E)$.
Of course we have $E = \Hom_R(R, E) = D(R)$. Since $E$ is injective
the functor $D$ is exact. The result now follows immediately from the
description of the categories in
Lemma \ref{lemma-describe-categories}.
\end{proof}
\end{document}
